\section{PECL as a Constrained Update Geometry}
\label{sec:pecl-geometry}

LoRA-based continual learning turns sequential adaptation into a constrained geometric problem.
Unlike full-parameter fine-tuning, future tasks do not modify the model freely in the ambient parameter space, but only through low-rank update directions induced by the trainable LoRA parameters.
To study generalization and forgetting in this regime, we therefore need to make explicit not only which parameters are trainable, but also which \emph{weight-space directions} are actually realizable under the PECL mechanism.

\subsection{Frozen-Backbone PECL as Constrained Adaptation}
\label{sec:decomp}

Let $\mathcal{X}$, $\mathcal{Y}$, and $\mathcal{Z} = \mathcal{X} \times \mathcal{Y}$ denote the input, label, and data spaces.
A model with parameters $W \in \mathbb{R}^d$ induces a predictor $f_W : \mathcal{X} \to \mathcal{Y}$.
We fix a bounded loss $\ell : \mathcal{H} \times \mathcal{Z} \to [0,1]$ and write $\ell(W,z)$ for $\ell(f_W,z)$.
In continual learning with $T$ tasks, each task $t$ provides a sample $S_t = \{z_{tj}\}_{j=1}^{m_t}$ drawn i.i.d.\ from distribution $\mathcal{D}_t$.

In frozen-backbone PECL, the model at task $t$ decomposes as
\[
W_t = W_t^{\mathrm{froz}} + \Delta W_t,
\]
where $W_t^{\mathrm{froz}}$ collects all parameters frozen while learning task $t$, and $\Delta W_t$ is the trainable increment produced by the current-task LoRA adapter.
The key PECL constraint is that the update is not arbitrary.
Rather, it is restricted to a task-dependent admissible update geometry:
\[
\Delta W_t \in \mathcal{W}_{\Delta,t} \subseteq \mathbb{R}^{d_1 \times d_2},
\]
where $\mathcal{W}_{\Delta,t}$ is a linear subspace in the ambient weight-increment space.
Accordingly, the hypothesis class available at task $t$ is the affine set $W_t^{\mathrm{froz}} + \mathcal{W}_{\Delta,t}$, and all task-dependent adaptation is confined to $\mathcal{W}_{\Delta,t}$.

This is the central geometric fact underlying the rest of the paper.
Under the frozen-backbone constraint, future tasks can only adapt the model through perturbations supported on $\mathcal{W}_{\Delta,t}$.
\emph{Any curvature or flatness property aggregated over directions outside $\mathcal{W}_{\Delta,t}$ is not automatically exploitable by future continual adaptation.}

\subsection{From LoRA Coordinates to Admissible Update Geometry}
\label{sec:lora-bridge}

The central object of our analysis is not the LoRA factor-coordinate space itself, but the set of \emph{weight-increment directions} that the LoRA parameterization can realize locally around the learned solution.

Let $\theta_{\Delta} = (A,B)$ denote the trainable LoRA coordinates, and let the corresponding weight increment be
\[
\Delta W(\theta_{\Delta}) = AB.
\]
Since the map $\theta_{\Delta} \mapsto \Delta W(\theta_{\Delta})$ is bilinear, the globally reachable set is the variety of rank-$\le r$ matrices---not a linear subspace.
For our PAC-Bayesian analysis, however, we only require a \emph{local linear model} of the admissible update directions around the learned task-$t$ solution.

Let $\hat{\theta}_{\Delta,t} = (\hat A_t, \hat B_t)$ be the learned LoRA parameters after task $t$.
A first-order expansion yields
\[
\Delta W(\hat{\theta}_{\Delta,t} + \delta\theta_{\Delta})
= \Delta W(\hat{\theta}_{\Delta,t})
+ J_{\mathrm{LoRA}}(\hat{\theta}_{\Delta,t})\,\delta\theta_{\Delta}
+ O(\|\delta\theta_{\Delta}\|^2),
\]
where $J_{\mathrm{LoRA}}(\hat{\theta}_{\Delta,t}) = [\hat B_t^\top \otimes I_{d_1},\; I_{d_2} \otimes \hat A_t]$ is the Jacobian of the parameterization map at $\hat\theta_{\Delta,t}$.
This motivates defining the task-admissible update subspace as
\[
\mathcal{W}_{\Delta,t}
:= \mathrm{Im}\!\left(J_{\mathrm{LoRA}}(\hat{\theta}_{\Delta,t})\right)
\subseteq \mathbb{R}^{d_1 \times d_2}.
\]
Thus, $\mathcal{W}_{\Delta,t}$ is the local tangent subspace in weight-increment space induced by the LoRA parameterization at the learned solution.
It has dimension at most $r(d_1 + d_2)$, is fully determined by $(\hat A_t, \hat B_t)$, and changes after each task.

This distinction is crucial for the rest of the paper.
The PAC-Bayesian theory is stated on $\mathcal{W}_{\Delta,t}$, because the bound concerns perturbations and posterior complexity in the space of realizable weight increments.
In implementation, perturbations are applied in the trainable LoRA coordinates $(A,B)$.
Under the local linearization above, such coordinate perturbations induce increment perturbations supported on $\mathcal{W}_{\Delta,t}$---so trainable-coordinate perturbations admit a local geometric interpretation through $\mathcal{W}_{\Delta,t}$.
However, trainable-coordinate flatness and admissible-update flatness should not be conflated as globally identical: the former is an implementation-level quantity defined in the optimization coordinate space, while the latter is the geometric object studied by our theory, defined in weight-increment space and invariant to local reparameterization.
See Appendix~\ref{app:lora_subspace_bridge} for the full derivation.

\textbf{Theory object vs.\ empirical proxy.}
Our PAC-Bayesian analysis is stated on the abstract $\mathcal{W}_{\Delta,t}$, while empirical diagnostics such as the amplification factor $\alpha_t$ and curvature projection ratio $c_t$ are constructed from the learned $\Delta W_t$ as an empirical proxy for $\mathcal{W}_{\Delta,t}$.
These proxies are justified by the Jacobian bridge, but the distinction should be kept in mind when interpreting the experimental results in Section~\ref{sec:experiments}.

\subsection{Geometry Evolution Across Tasks}
\label{sec:lora-orgs}

Different LoRA-based continual learning methods can be understood as different rules for how the admissible update geometry $\mathcal{W}_{\Delta,t}$ evolves across tasks.

\textbf{SeqLoRA (fixed geometry).}
A single LoRA module is updated sequentially across all tasks while the backbone remains frozen, yielding $W = W_0 + AB$.
The admissible geometry is heavily shared across tasks, so cross-task overlap is maximal and new tasks are most likely to reuse or overwrite previously occupied directions.

\textbf{IncLoRA (expanding geometry).}
A new LoRA module is added for each task while previous modules and the backbone are frozen, yielding $W = W_0 + \sum_{t=1}^T A_t B_t$.
The admissible geometry grows over time, allowing new tasks to access additional update directions and reducing direct overlap between task-specific updates.

\textbf{Constrained expansion.}
A third family retains the incremental design of IncLoRA but imposes structural constraints on $\mathcal{W}_{\Delta,t}$ to limit cross-task interference.
The constraint mechanism differs across methods:
\begin{itemize}
  \item \emph{OLoRA}~\citep{wang2023orthogonal} acts in parameter space, penalizing overlap between adapter directions via $\sum_{i < t}\|A_i^\top A_t\|_F^2$ so that newly added directions are disentangled from previous ones.
  \item \emph{InfLoRA}~\citep{Liang_2024_CVPR} acts in gradient space, projecting each task's gradient into the null space of the subspace spanned by previous task gradients.
  \item \emph{SDLoRA}~\citep{wu2025sdlorascalabledecoupledlowrank} acts through component decomposition, splitting the LoRA update into task-shared and task-specific parts with independently controlled ranks.
\end{itemize}

Although these methods differ algorithmically, they can all be understood as defining different evolution rules for the sequence $\{\mathcal{W}_{\Delta,t}\}_{t=1}^T$.
As we show in Section~\ref{sec:pecl-spec}, the PAC-Bayesian reduction to $\mathcal{W}_{\Delta,t}$ holds for all three organizations; the organization determines the structure of $\mathcal{W}_{\Delta,t}$ itself, and hence how sharply the scope-type distinction manifests empirically.
